% GENERATED FILE. Edit canonical lessons, then run npm run generate:books.
% canonical-source-hash: fnv1a64:a5d989c0ce236d95
% canonical-lessons: MW-C09-hear-nana, MW-C09-nana, MW-C09-hear-nani, MW-C09-nani, MW-C09-hear-pati, MW-W09-ta, MW-C09-pati, MW-R09-maternal-three

\chapter{Maternal Grandparents and Husband: One New Sign}
\label{ch:mw-family-ten}
\hlchaptermodality{12}

\begin{chapteropening}
\textbf{By the end of this chapter:} \emph{I can add maternal grandfather, maternal grandmother, and husband to the family map after learning only one new Devanagari sign.}

Three source-attested relationship labels begin by ear, reuse the established script except for \mw{त}, and finish in four separately scored skills.
\end{chapteropening}

\section[nānā]{Hear \emph{nānā} — maternal grandfather}

\label{lesson:MW-C09-hear-nana}



\begin{quote}
Say the paternal-grandfather word. \emph{Write it:} \textbf{\mw{परिवार}}

Then recall the seven-word family map. Today one matching two-beat word changes only the family side.
\end{quote}

\subsection*{What to know first}
\begin{quote}
\emph{nā-nā} — \textbf{maternal grandfather}
\end{quote}

Listen twice. Say it once. Keep \emph{dādā} and \emph{nānā} tied to their meanings before writing enters.

\begin{tcolorbox}[breakable,colback=teal!4,colframe=teal!35!black,title={Before you move on}]
Say \emph{nānā} once without seeing its spelling.

Source: \href{https://www.marwaripathshala.com/marwari-lesson-4-english}{Marwari Pathshala, Lesson 4}.
\end{tcolorbox}
\section[\texorpdfstring{\mw{नाना} (nānā)}{नाना (nānā)}]{\mw{नाना} — one known unit twice}

\label{lesson:MW-C09-nana}



\begin{quote}
Say \emph{nānā}. \emph{Write it:} familiar \textbf{\mw{न}} and \textbf{\mw{ा}}
\end{quote}

\subsection*{What to know first}
\begin{quote}
\textbf{\mw{ना}} + \textbf{\mw{ना}} $\to$ \textbf{\mw{नाना}} — \emph{nānā} — maternal grandfather
\end{quote}

\subsection*{Your turn}
\emph{Cover:} the word after a five-second look

\emph{Write it:} the word

Then label \textbf{\mw{दादा}} and \textbf{\mw{नाना}} orally without ranking either side of the family.

\begin{tcolorbox}[breakable,colback=teal!4,colframe=teal!35!black,title={Before you move on}]
Source: \href{https://www.marwaripathshala.com/marwari-lesson-4-english}{Marwari Pathshala, Lesson 4}.
\end{tcolorbox}
\section[nānī]{Hear \emph{nānī} — maternal grandmother}

\label{lesson:MW-C09-hear-nani}



\begin{quote}
Say \emph{nānā}. \emph{Write it:} \textbf{\mw{बाप}}

Then say the known paternal-grandmother word. Listen for the final long \emph{ī} in the new word.
\end{quote}

\subsection*{What to know first}
\begin{quote}
\emph{nānī} — \textbf{maternal grandmother}
\end{quote}

\begin{tcolorbox}[breakable,colback=teal!4,colframe=teal!35!black,title={Before you move on}]
Say \emph{nānī} once. Writing waits for the next step.

Source: \href{https://www.marwaripathshala.com/marwari-lesson-4-english}{Marwari Pathshala, Lesson 4}.
\end{tcolorbox}
\section[\texorpdfstring{\mw{नानी} (nānī)}{नानी (nānī)}]{\mw{नानी} — change only the final vowel}

\label{lesson:MW-C09-nani}



\begin{quote}
\emph{Write it:} \textbf{\mw{नाना}}, then familiar \textbf{\mw{ी}} separately

Then recall the seven-label four-skill payoff.
\end{quote}

\subsection*{What to know first}
\begin{quote}
\textbf{\mw{ना}} + \textbf{\mw{नी}} $\to$ \textbf{\mw{नानी}} — \emph{nānī} — maternal grandmother
\end{quote}

\subsection*{Your turn}
\emph{Cover:} \textbf{\mw{नानी}}, then wait five seconds

\emph{Write it:} the word, and circle only the final \textbf{\mw{ी}}

\begin{tcolorbox}[breakable,colback=teal!4,colframe=teal!35!black,title={Before you move on}]
Source: \href{https://www.marwaripathshala.com/marwari-lesson-4-english}{Marwari Pathshala, Lesson 4}.
\end{tcolorbox}
\section[pati]{Hear \emph{pati} — husband}

\label{lesson:MW-C09-hear-pati}



\begin{quote}
Say the two maternal-grandparent labels. \emph{Write it:} \textbf{\mw{नानी}}, then familiar \textbf{\mw{र}} and \textbf{\mw{भाई}} from sound cues
\end{quote}

\subsection*{What to know first}
\begin{quote}
\emph{pati} — \textbf{husband}
\end{quote}

This is a relationship label, not a required household pattern. Say the two short beats. The spelling waits until its new consonant has its own lesson.

\begin{tcolorbox}[breakable,colback=teal!4,colframe=teal!35!black,title={Before you move on}]
Source: \href{https://www.marwaripathshala.com/marwari-lesson-4-english}{Marwari Pathshala, Lesson 4}.
\end{tcolorbox}
\section[\texorpdfstring{\mw{त} (ta)}{त (ta)}]{\mw{त} — a quiet dental stop}

\label{lesson:MW-W09-ta}



\begin{quote}
Say \emph{pati}. Write known \textbf{\mw{ठ}} and \textbf{\mw{थ}} once; today's unaspirated dental sound gets its own sign.
\end{quote}

\subsection*{Script — one sign}
\hlblockfigure{\detokenize{figures/MW-W09-ta-filmstrip.pdf}}{How \mw{त} is written, stroke by stroke}

\begin{quote}
\textbf{\mw{त}} — \emph{ta}
\end{quote}

\subsection*{Your turn}
Trace once with the model visible, copy once, cover it, and write it once.

\begin{tcolorbox}[breakable,colback=teal!4,colframe=teal!35!black,title={Before you move on}]
Source: \href{https://www.unicode.org/charts/PDF/U0900.pdf}{Unicode Devanagari chart}.
\end{tcolorbox}
\section[\texorpdfstring{\mw{पति} (pati)}{पति (pati)}]{\mw{पति} — husband on the page}

\label{lesson:MW-C09-pati}



\begin{quote}
Say husband and maternal grandfather. \emph{Write it:} \textbf{\mw{प}}, \textbf{\mw{त}}, \textbf{\mw{ि}}, \textbf{\mw{राम}}, and \textbf{\mw{बहन}} — point to \textbf{\mw{म}} inside \textbf{\mw{राम}}
\end{quote}

\subsection*{What to know first}
\begin{quote}
\textbf{\mw{प}} + \textbf{\mw{ति}} $\to$ \textbf{\mw{पति}} — \emph{pati} — husband
\end{quote}

\subsection*{Your turn}
\emph{Cover:} the model after one look, then wait five seconds

\emph{Write it:} \textbf{\mw{पति}} — check that \textbf{\mw{ि}} appears before \textbf{\mw{त}} on the page but is read after it

\begin{tcolorbox}[breakable,colback=teal!4,colframe=teal!35!black,title={Before you move on}]
Source: \href{https://www.marwaripathshala.com/marwari-lesson-4-english}{Marwari Pathshala, Lesson 4}.
\end{tcolorbox}
\section[Practice]{Checkpoint — three more relationship labels}

\label{lesson:MW-R09-maternal-three}



\begin{quote}
Say mother and father. \emph{Write it:} both words, then \textbf{\mw{स}}

Recall the earlier four-label payoff before extending the map.
\end{quote}

\subsection*{Your turn}
\begin{enumerate}
\raggedright
  \item Hear and identify all three labels.
  \item Produce each from a meaning cue.
  \item \emph{Read it:} \textbf{\mw{नाना} — \mw{नानी} — \mw{पति}}
  \item \emph{Cover:} the line
  \item \emph{Write it:} all three
\end{enumerate}

\begin{tcolorbox}[breakable,colback=teal!4,colframe=teal!35!black,title={Before you move on}]
Source: \href{https://www.marwaripathshala.com/marwari-lesson-4-english}{Marwari Pathshala, Lesson 4}.
\end{tcolorbox}
